\documentclass[8pt]{extarticle}

% --- Math + typography
\usepackage{amsmath,amssymb}
\usepackage[margin=0.40in]{geometry}
\usepackage{microtype}
\linespread{0.95}
\setlength{\parskip}{1.8pt}
\setlength{\parindent}{0pt}

% --- Compact display math, arrays, aligns
\makeatletter
\g@addto@macro\normalsize{%
  \setlength\abovedisplayskip{3pt}%
  \setlength\belowdisplayskip{3pt}%
  \setlength\abovedisplayshortskip{2pt}%
  \setlength\belowdisplayshortskip{2pt}%
}
\makeatother
\setlength{\arraycolsep}{2pt}
\renewcommand{\arraystretch}{0.93}
\setlength{\jot}{1.5pt}
\mathsurround=0pt

% --- Lists tighter
\usepackage{enumitem}
\setlist{nosep,leftmargin=*}

% --- Boxes + columns
\usepackage[most]{tcolorbox}
\usepackage{multicol}
\setlength{\columnsep}{0.22em}

\tcbset{
  colback=white,
  colframe=black!55,
  boxrule=0.35pt,
  arc=0.8pt,
  boxsep=1.2pt,
  left=1.2pt,right=1.2pt,top=1.2pt,bottom=1.2pt,
  before skip=2.2pt, after skip=2.2pt,
  enhanced,
  breakable,
  fonttitle=\bfseries\footnotesize,
  title style={opacityfill=0}
}

\begin{document}
\scriptsize

\begin{center}
  \normalsize\textbf{Math 262 -- Solutions to Practice Problems}
\end{center}

\begin{multicols}{2}
\raggedcolumns

% ---------------------- Problem 1 ----------------------
\begin{tcolorbox}[title=]
\noindent
\textbf{1.} Let \( T:\mathbb{R}^3 \to \mathbb{R}^2 \) be given by
\[
T\!\left(\begin{bmatrix}x_{1}\\x_{2}\\x_{3}\end{bmatrix}\right)
=\begin{bmatrix}x_{1}+x_{2}\\ x_{2}+x_{3}\end{bmatrix}.
\]
\begin{enumerate}
\item[(a)] Show that \(T\) is a linear transformation.
\item[(b)] Find the matrix of \(T\).
\end{enumerate}

\textbf{Solution.}
\begin{enumerate}
\item
\[
a^{(*)}\;T\!\left(\begin{bmatrix}x_{1}\\x_{2}\\x_{3}\end{bmatrix}
+\begin{bmatrix}y_{1}\\y_{2}\\y_{3}\end{bmatrix}\right)
=T\!\left(\begin{bmatrix}x_{1}+y_{1}\\x_{2}+y_{2}\\x_{3}+y_{3}\end{bmatrix}\right)
\]
\[
=\begin{bmatrix}x_{1}+y_{2}\\x_{2}+y_{3}\end{bmatrix}
-\begin{bmatrix}(x_{1}+y_{1})+(x_{2}+y_{2})\\(x_{2}+y_{2})+(x_{3}+y_{3})\end{bmatrix}
\]
\[
=\begin{bmatrix}x_{1}+y_{2}\\x_{2}+y_{3}\end{bmatrix}
-\begin{bmatrix}x_{1}+x_{2}+y_{1}+y_{2}\\x_{2}+x_{3}+y_{2}+y_{3}\end{bmatrix}
= T\!\left(\begin{bmatrix}x_{1}\\x_{2}\\x_{3}\end{bmatrix}\right)
+T\!\left(\begin{bmatrix}y_{1}\\y_{2}\\y_{3}\end{bmatrix}\right).
\]
(i) For \(r\in\mathbb{R}\):
\[
T\!\left(r\begin{bmatrix}x_{1}\\x_{2}\\x_{3}\end{bmatrix}\right)
=\begin{bmatrix}rx_{1}+rx_{2}\\rx_{2}+rx_{3}\end{bmatrix}
-r\begin{bmatrix}x_{1}+x_{2}\\x_{2}+x_{3}\end{bmatrix}
=r\,T\!\left(\begin{bmatrix}x_{1}\\x_{2}\\x_{3}\end{bmatrix}\right).
\]
So \(T\) is linear.
\item
\[
T(e_{1})=\begin{bmatrix}1\\0\end{bmatrix},\quad
T(e_{2})=\begin{bmatrix}1\\1\end{bmatrix},\quad
T(e_{3})=\begin{bmatrix}0\\1\end{bmatrix}
\Rightarrow
A=\begin{bmatrix}1&1&0\\0&1&1\end{bmatrix}.
\]
\end{enumerate}
\end{tcolorbox}

% ---------------------- Problem 2 ----------------------
\begin{tcolorbox}[title=]
\noindent
\textbf{2.} Given \(A\vec{x}=\vec{b}\), suppose the rref of the augmented matrix is
\[
\left[
\begin{array}{cccc|c}
1&0&0&1&0\\
0&1&0&2&1\\
0&0&1&3&2\\
0&0&0&0&0
\end{array}
\right].
\]
Give the general solution to the system.

\textbf{Solution.}
\[
\begin{cases}
x_{1}+x_{4}=0\\
x_{2}+2x_{4}=1\\
x_{3}+3x_{4}=2
\end{cases}
\Rightarrow
\begin{aligned}
x_{1}&=-x_{4}\\
x_{2}&=2x_{4}+1\\
x_{3}&=-3x_{4}+2
\end{aligned}
\quad(x_{4}=r).
\]
\[
\begin{bmatrix}x_{1}\\x_{2}\\x_{3}\\x_{4}\end{bmatrix}
=\begin{bmatrix}-r\\2r+1\\-3r+2\\r\end{bmatrix}.
\]
\end{tcolorbox}

% ---------------------- Problem 3 (full, with row ops) ----------------------
\begin{tcolorbox}[title=]
\noindent
\textbf{Problem 3 (Original Statement):}
Given
\(
\begin{cases}
x_{1}-x_{3}=1,\\
2x_{1}-2x_{2}=1,\\
3x_{1}-4x_{2}+x_{3}=a,
\end{cases}
\)
find (a) no solution, (b) exactly one solution, (c) infinitely many solutions; give the general solution when infinite.

\textbf{Solution.}
\[
\left[
\begin{array}{ccc|c}
1&0&-1&1\\
-2&2&0&1\\
3&-4&1&a
\end{array}
\right]
\]
Row ops: \(R_{2}\to R_{2}+2R_{1}\), \(R_{3}\to R_{3}-3R_{1}\).
\[
\left[
\begin{array}{ccc|c}
1&0&-1&1\\
0&-2&2&-1\\
0&-4&4&a-3
\end{array}
\right]
\]
Row op: \(R_{2}\to (-\tfrac12)R_{2}\).
\[
\left[
\begin{array}{ccc|c}
1&0&-1&1\\
0&1&-1&\tfrac12\\
0&-4&4&a-3
\end{array}
\right]
\]
Row op: \(R_{3}\to R_{3}+4R_{2}\).
\[
\left[
\begin{array}{ccc|c}
1&0&-1&1\\
0&1&-1&\tfrac12\\
0&0&0&a-1
\end{array}
\right]
\]
(a) No solution if \(a\neq1\). \quad (b) Never exactly one. \quad (c) Infinitely many if \(a=1\).

If \(a=1\), let \(x_{3}=r\). Then \(x_{2}=\tfrac12+r\), \(x_{1}=1+r\), so
\(
\begin{bmatrix}x_{1}\\x_{2}\\x_{3}\end{bmatrix}
=\begin{bmatrix}1+r\\ \tfrac12+r\\ r\end{bmatrix}.
\)
\end{tcolorbox}

% ---------------------- Problem 4 ----------------------
\begin{tcolorbox}[title=]
\noindent
In \(\mathbb{R}^{2}\), let
\(
\vec{u}=\begin{bmatrix}\tfrac{1}{\sqrt{10}}\\[4pt]\tfrac{3}{\sqrt{10}}\end{bmatrix}.
\)
\begin{enumerate}
\item[(a)] Show that \(\vec{u}\) is a unit vector.
\item[(b)] Let \(L\) be the line spanned by \(\vec{u}\), and let
\(
\vec{x}=\begin{bmatrix}-1\\4\end{bmatrix}.
\)
Find the reflection \(\mathrm{ref}_{L}(\vec{x})\).
\end{enumerate}

\textbf{Solution.}
\[
|\vec{u}|=\sqrt{\big(\tfrac1{\sqrt{10}}\big)^{2}+\big(\tfrac3{\sqrt{10}}\big)^{2}}=1,\qquad
\mathrm{proj}_{L}(\vec{x})=\frac{\vec{x}\cdot\vec{u}}{\vec{u}\cdot\vec{u}}\vec{u}.
\]
\[
\vec{x}\cdot\vec{u}
=\begin{bmatrix}-1\\4\end{bmatrix}\!\cdot\!
\begin{bmatrix}\tfrac1{\sqrt{10}}\\[2pt]\tfrac3{\sqrt{10}}\end{bmatrix}
=-\tfrac1{\sqrt{10}}+\tfrac{12}{\sqrt{10}}=\tfrac{11}{\sqrt{10}},
\quad
\mathrm{proj}_{L}(\vec{x})
=\begin{bmatrix}\tfrac{11}{10}\\[2pt]\tfrac{33}{10}\end{bmatrix}.
\]
\[
\mathrm{ref}_{L}(\vec{x})=2\,\mathrm{proj}_{L}(\vec{x})-\vec{x}.
\]
\end{tcolorbox}

% ---------------------- Problem 5 (full original content) ----------------------
\begin{tcolorbox}[title=]
\noindent
\textbf{Original Problem:}
In \(\mathbb{R}^{3}\), let \(L\) be the line spanned by
\(
\vec{w}=\begin{bmatrix}\tfrac1{\sqrt{2}}\\[2pt]0\\[2pt]\tfrac1{\sqrt{2}}\end{bmatrix}
\),
\(
\vec{x}=\begin{bmatrix}1\\1\\1\end{bmatrix}.
\)
Verify that \(\vec{x}-\mathrm{proj}_{L}(\vec{x})\) is orthogonal to \(\vec{w}\).

\textbf{Solution.}
\[
2\begin{bmatrix}\frac{11}{10}\\[2pt]\frac{33}{10}\end{bmatrix}
-
\begin{bmatrix}-1\\[2pt]4\end{bmatrix}
=
\begin{bmatrix}\frac{22}{10}\\[2pt]\frac{66}{10}\end{bmatrix}
-
\begin{bmatrix}-1\\[2pt]4\end{bmatrix}
=
\begin{bmatrix}\frac{22}{10}+1\\[2pt]\frac{66}{10}-4\end{bmatrix}
=
\begin{bmatrix}\frac{16}{5}\\[2pt]\frac{13}{5}\end{bmatrix}.
\]
\textit{Note: You could also use the matrix formula for this one.}

\[
\mathrm{proj}_{L}(\vec{x})
=\frac{\vec{x}\cdot\vec{w}}{\vec{w}\cdot\vec{w}}\vec{w}
= \frac{2}{\sqrt{2}}
\begin{bmatrix}\frac{1}{\sqrt{2}}\\[2pt]\frac{1}{\sqrt{2}}\end{bmatrix}
= \begin{bmatrix}1\\1\end{bmatrix}.
\]
\[
\vec{x}-\mathrm{proj}_{L}(\vec{x})
=\begin{bmatrix}1\\1\\1\end{bmatrix}
-\begin{bmatrix}0\\1\\0\end{bmatrix}
=\begin{bmatrix}1\\0\\1\end{bmatrix}.
\]
\[
\begin{bmatrix}1\\0\\1\end{bmatrix}\cdot
\begin{bmatrix}\tfrac1{\sqrt{2}}\\[2pt]0\\[2pt]\tfrac1{\sqrt{2}}\end{bmatrix}=0,
\quad \text{so it is orthogonal.}
\]
\end{tcolorbox}

% ---------------------- Problem 6 ----------------------
\begin{tcolorbox}[title=]
\noindent
\textbf{Problem 6 (Original Statement):}
Define \(T:\mathbb{R}^{2}\to\mathbb{R}^{2}\) by
\(
T\big(\begin{bmatrix}x_{1}\\x_{2}\end{bmatrix}\big)
=\begin{bmatrix}x_{1}-x_{2}\\x_{1}+x_{2}\end{bmatrix}.
\)
Assume \(T\) is linear.
\begin{enumerate}
\item[(a)] Find \(A\), the matrix of \(T\).
\item[(b)] Use \(A\) to explain why \(T\) is invertible.
\item[(c)] Use \(A\) to find the rule for \(T^{-1}\).
\end{enumerate}

\textbf{Solution.}
\[
T(e_{1})=\begin{bmatrix}1\\1\end{bmatrix},\quad
T(e_{2})=\begin{bmatrix}-1\\1\end{bmatrix},\quad
A=\begin{bmatrix}1&-1\\1&1\end{bmatrix}.
\]
\[
\det A=(1)(1)-(-1)(1)=2\Rightarrow A\text{ invertible},\;
T^{-1}=T_{A^{-1}},\;
A^{-1}=\tfrac12\begin{bmatrix}1&1\\-1&1\end{bmatrix}.
\]
\[
T^{-1}\!\begin{bmatrix}x_{1}\\x_{2}\end{bmatrix}
=\begin{bmatrix}\tfrac12(x_{1}+x_{2})\\[2pt]\tfrac12(-x_{1}+x_{2})\end{bmatrix}.
\]
\end{tcolorbox}

% ---------------------- Problem 7 (restored full original solution) ----------------------
\begin{tcolorbox}[title=]
\noindent
\textbf{Problem 7 (Original Statement):}  

Let \(\theta\) be any angle and let \(A\) be the matrix which represents the counterclockwise rotation through \(\theta\).  
Prove that for any vector
\[
\vec{u} =
\begin{bmatrix} x\\ y \end{bmatrix} \in \mathbb{R}^{2},
\quad
|A\vec{u}| = |\vec{u}|
\]
—that is, rotations preserve length.

\bigskip
\noindent
\textbf{Solution:}

\[
A =
\begin{bmatrix}
\cos\theta & -\sin\theta\\
\sin\theta & \cos\theta
\end{bmatrix},
\qquad
\vec{u} =
\begin{bmatrix} x\\ y \end{bmatrix},
\qquad
A\vec{u} =
\begin{bmatrix}
x\cos\theta - y\sin\theta\\
x\sin\theta + y\cos\theta
\end{bmatrix}.
\]

\[
|A\vec{u}|
= \sqrt{(x\cos\theta - y\sin\theta)^{2} + (x\sin\theta + y\cos\theta)^{2}}
\]

Expand:
\[
= \sqrt{
x^{2}\cos^{2}\theta - 2xy\cos\theta\sin\theta + y^{2}\sin^{2}\theta
+ x^{2}\sin^{2}\theta + 2xy\sin\theta\cos\theta + y^{2}\cos^{2}\theta
}
\]

\[
= \sqrt{
x^{2}(\cos^{2}\theta+\sin^{2}\theta) + y^{2}(\sin^{2}\theta+\cos^{2}\theta)
}
= \sqrt{x^{2} + y^{2}}
= |\vec{u}|.
\]

Hence rotations preserve length.
\end{tcolorbox}
% ---------------------- Problem 8 ----------------------
\begin{tcolorbox}[title=]
\noindent
\textbf{Problem 8 (Original Statement):}
Fill in the blanks for an invertible \(n\times n\) matrix \(A\):
\begin{enumerate}
\item[(a)] \(\mathrm{rank}(A)=\underline{\hspace{1cm}}\)
\item[(b)] The homogeneous system \(A\vec{x}=\vec{0}\) has only the \underline{\hspace{2cm}} solution.
\item[(c)] \(A\) is row-reducible to the \underline{\hspace{2cm}} matrix.
\end{enumerate}

\textbf{Solution.}
\[
\mathrm{rank}(A)=n,\qquad \text{only the trivial solution},\qquad
A\text{ row-reduces to the identity}.
\]
\end{tcolorbox}

% ---------------------- Problem 9 (with row ops) ----------------------
\begin{tcolorbox}[title=]
\noindent
Let
\(
A=\begin{bmatrix}
1&1&0&1&0\\
2&3&4&0&1\\
3&4&0&1&5\\
6&8&4&2&6
\end{bmatrix}.
\)
Give the reduced row echelon form of \(A\).

\textbf{Solution.}
\[
\begin{bmatrix}
1&1&0&1&0\\
2&3&4&0&1\\
3&4&0&1&5\\
6&8&4&2&6
\end{bmatrix}
\]
Row ops: \(R_{2}\!\to\!R_{2}-2R_{1}\), \(R_{3}\!\to\!R_{3}-3R_{1}\), \(R_{4}\!\to\!R_{4}-6R_{1}\).
\[
\begin{bmatrix}
1&1&0&1&0\\
0&1&4&-2&1\\
0&1&0&-2&5\\
0&2&4&-4&6
\end{bmatrix}
\]
Row ops: \(R_{3}\!\to\!R_{3}-R_{2}\), \(R_{4}\!\to\!R_{4}-2R_{2}\).
\[
\begin{bmatrix}
1&1&0&1&0\\
0&1&4&-2&1\\
0&0&-4&0&4\\
0&0&-4&0&4
\end{bmatrix}
\]
Row ops: \(R_{3}\!\to\!(-\tfrac14)R_{3}\), \(R_{4}\!\to\!R_{4}-R_{3}\).
\[
\begin{bmatrix}
1&1&0&1&0\\
0&1&4&-2&1\\
0&0&1&0&-1\\
0&0&0&0&0
\end{bmatrix}
\]
Row op: \(R_{1}\!\to\!R_{1}-4R_{3}\).
\[
\boxed{
\begin{bmatrix}
1&0&0&-5\\
0&1&0&5\\
0&0&1&-1\\
0&0&0&0
\end{bmatrix}}
\]
\end{tcolorbox}

\begin{tcolorbox}[title=Quick Determinant Tips, colframe=black!60!white, colback=gray!2]
\small
\(\det(A)\) changes sign when rows are swapped;  
scaling a row by \(k\) multiplies \(\det\) by \(k\);  
if two rows are identical or one is all zeros → \(\det(A)=0\).
\end{tcolorbox}

% ---------------------- Bonus: Nonlinear Example ----------------------
\begin{tcolorbox}[title=Example – Transformation That Is Not Linear]
Let \(T:\mathbb{R}^2\to\mathbb{R}\) be defined by \(T(x_1,x_2)=x_1^2+x_2.\)
\[
T(\vec{x}+\vec{y})=(x_1+y_1)^2+(x_2+y_2)
=x_1^2+2x_1y_1+y_1^2+x_2+y_2
\neq T(\vec{x})+T(\vec{y})
\]
(because of the \(2x_1y_1\) term).  
Thus \(T\) fails additivity ⇒ not linear.
\end{tcolorbox}

% ---------------------- Bonus: Subspace Proof ----------------------
\begin{tcolorbox}[title=Proof – Orthogonal Set Is a Subspace]
Let \(\vec u\in\mathbb{R}^2\) and define \(S=\{\vec x\in\mathbb{R}^2\mid \vec x\cdot\vec u=0\}\).
\begin{itemize}
\item If \(\vec x,\vec y\in S\), then  
\((\vec x+\vec y)\cdot\vec u
= \vec x\cdot\vec u+\vec y\cdot\vec u=0+0=0\)
⇒ \(\vec x+\vec y\in S.\)
\item If \(k\in\mathbb{R}\) and \(\vec x\in S\), then  
\((k\vec x)\cdot\vec u
= k(\vec x\cdot\vec u)=k(0)=0\)
⇒ \(k\vec x\in S.\)
\end{itemize}
Hence \(S\) is closed under addition and scalar multiplication ⇒ \(S\) is a subspace.
\end{tcolorbox}


% ---------------------- Problem 10(a) (1:1 handwritten version) ----------------------
% ---------------------- Problem 10(a,b) — Original Statements and Solutions (compressed full) ----------------------
% ---------------------- Problem 10(a,b) — Original Statements and Soluti% ---------------------- Problem 10(a,b) — Original Statements and Solutions (compact, matches original) ----------------------
\begin{tcolorbox}[title=]
\noindent
\textbf{Problem 10(a) (Original Statement):}\quad
Let
\[
A=\begin{bmatrix}1&2\\[2pt]2&1\end{bmatrix}.
\]
Prove that if
\[
AB=\begin{bmatrix}0&0\\[2pt]0&0\end{bmatrix},
\]
then
\[
B=\begin{bmatrix}0&0\\[2pt]0&0\end{bmatrix}.
\]

\smallskip
\textbf{Solution.}
Let \(B=\begin{bmatrix}a&b\\ c&d\end{bmatrix}\).
Then
\[
AB=
\begin{bmatrix}1&2\\ 2&1\end{bmatrix}
\begin{bmatrix}a&b\\ c&d\end{bmatrix}
=
\begin{bmatrix}a+2c & b+2d\\[2pt] 2a+c & 2b+d\end{bmatrix}.
\]
Setting \(AB=0\) gives
\[
a+2c=0,\;2a+c=0 \ \Rightarrow\ a=c=0, \qquad
b+2d=0,\;2b+d=0 \ \Rightarrow\ b=d=0,
\]
so \(B=0\).

\medskip
\textbf{Problem 10(b) (Original Statement):}\quad
Let
\[
C=\begin{bmatrix}0&1\\[2pt]0&1\end{bmatrix}.
\]
Give the general form of a matrix \(D\) such that
\[
\boxed{\,C D=\begin{bmatrix}0&0\\[2pt]0&0\end{bmatrix}\,}.
\]

\smallskip
\textbf{Solution.}
Let \(D=\begin{bmatrix}a&b\\ c&d\end{bmatrix}\).
Then
\[
C D=
\begin{bmatrix}0&1\\ 0&1\end{bmatrix}
\begin{bmatrix}a&b\\ c&d\end{bmatrix}
=
\begin{bmatrix}c&d\\ c&d\end{bmatrix}.
\]
For \(C D=0\) we need \(c=0\) and \(d=0\).
Hence the general form is
\[
D=\begin{bmatrix}a&b\\ 0&0\end{bmatrix}.
\]
\end{tcolorbox}
% ---------------------- Problem 11 (match midtermOnePage) ----------------------
\begin{tcolorbox}[title=]
\noindent
\textbf{Problem 11 (Original Statement):}  

In \(\mathbb{R}^{2}\), let
\[
\vec{u}=
\begin{bmatrix}1\\2\end{bmatrix},\qquad
\vec{v}=
\begin{bmatrix}2\\1\end{bmatrix}.
\]
Show that any vector
\[
\vec{w}=
\begin{bmatrix}a\\b\end{bmatrix}
\]
can be written as a linear combination
\(\vec{w}=r\vec{u}+s\vec{v}\).

\bigskip
\noindent
\textbf{Solution:}

We need to find scalars \(r\) and \(s\) such that
\[
r
\begin{bmatrix}1\\2\end{bmatrix}
+
s
\begin{bmatrix}2\\1\end{bmatrix}
=
\begin{bmatrix}a\\b\end{bmatrix}.
\]
This gives the system
\[
\begin{cases}
r+2s=a,\\
2r+s=b.
\end{cases}
\]

Multiply the first equation by \(2\):
\[
2r+4s=2a
\]
Subtract the second equation:
\[
(2r+4s)-(2r+s)=2a-b
\quad\Rightarrow\quad
3s=2a-b
\quad\Rightarrow\quad
s=\dfrac{2a-b}{3}.
\]

Substitute into \(r+2s=a\):
\[
r+2\!\left(\dfrac{2a-b}{3}\right)=a
\quad\Rightarrow\quad
3r+4a-2b=3a
\quad\Rightarrow\quad
3r=-a+2b
\quad\Rightarrow\quad
r=\dfrac{-a+2b}{3}.
\]

Hence
\[
\begin{bmatrix}a\\b\end{bmatrix}
=
\left(\dfrac{-a+2b}{3}\right)
\begin{bmatrix}1\\2\end{bmatrix}
+
\left(\dfrac{2a-b}{3}\right)
\begin{bmatrix}2\\1\end{bmatrix}.
\]
\end{tcolorbox}

% ---------------------- Problem 12 (full, with row ops & wrapped) ----------------------
\begin{tcolorbox}[title=]
\noindent
\begin{enumerate}
\item[(a)] Is
\(
\begin{bmatrix}
1&2&3\\
4&5&6\\
7&8&9
\end{bmatrix}
\)
invertible? Why or why not?
\item[(b)] Let
\(A=\begin{bmatrix}1&1&1\\2&1&0\\0&1&1\end{bmatrix}\).
Find \(A^{-1}\) and solve
\(A\vec{x}=\begin{bmatrix}1\\-3\\2\end{bmatrix}\).
\end{enumerate}

\textbf{Solution.}
\[
\left[
\begin{array}{ccc|ccc}
1&2&3&1&0&0\\
4&5&6&0&1&0\\
7&8&9&0&0&1
\end{array}
\right]
\]
Row ops: \(R_2\to R_2-4R_1\), \(R_3\to R_3-7R_1\).
\[
\left[
\begin{array}{ccc|ccc}
1&2&3&1&0&0\\
0&-3&-6&-4&1&0\\
0&-6&-12&-7&0&1
\end{array}
\right]
\]
Row op: \(R_3\to R_3-2R_2\).
\[
\left[
\begin{array}{ccc}
1&2&3\\
0&-3&-6\\
0&0&0
\end{array}
\right]
\Rightarrow \text{not invertible.}
\]

\[
\left[
\begin{array}{ccc|ccc}
1&1&1&1&0&0\\
2&1&0&0&1&0\\
0&1&1&0&0&1
\end{array}
\right]
\]
Row op: \(R_2\to R_2-2R_1\).
\[
\left[
\begin{array}{ccc|ccc}
1&1&1&1&0&0\\
0&-1&-2&-2&1&0\\
0&1&1&0&0&1
\end{array}
\right]
\]
Row ops: \(R_2\to -R_2\), \(R_1\to R_1-R_3\).
\[
\left[
\begin{array}{ccc|ccc}
1&0&-1&1&0&1\\
0&1&2&2&-1&0\\
0&1&1&0&0&1
\end{array}
\right]
\]
Row op: \(R_3\to R_3-R_2\).
\[
\left[
\begin{array}{ccc|ccc}
1&0&-1&1&0&1\\
0&1&2&2&-1&0\\
0&0&-1&-2&1&1
\end{array}
\right]
\]
Row op: \(R_3\to -R_3\).
\[
\left[
\begin{array}{ccc|ccc}
1&0&-1&1&0&1\\
0&1&2&2&-1&0\\
0&0&1&2&-1&-1
\end{array}
\right]
\]
Thus
\(A^{-1}=\begin{bmatrix}1&0&-1\\-2&1&2\\2&-1&-1\end{bmatrix}\),
\(
\vec{x}=A^{-1}\!\begin{bmatrix}1\\-3\\2\end{bmatrix}
=\begin{bmatrix}-1\\-1\\3\end{bmatrix}.
\)
\end{tcolorbox}

% ---------------------- Problem 13 ----------------------
\begin{tcolorbox}[title=]
\noindent
\textbf{Problem 13 (Original Statement):}
Let
\(A=\begin{bmatrix}1&r\\2r&1\end{bmatrix}\).
\begin{enumerate}
\item[(a)] Find \(\det A\).
\item[(b)] For which \(r\) is \(A\) invertible?
\item[(c)] When invertible, give \(A^{-1}\).
\end{enumerate}

\textbf{Solution.}
\[
\det A=(1)(1)-(2r)(r)=1-2r^{2}.
\]
\[
\det A=0\Rightarrow 1-2r^{2}=0\Rightarrow 2r^{2}=1\Rightarrow r=\pm\tfrac{1}{\sqrt{2}}.
\]
Thus \(A\) invertible for \(r\neq \pm\tfrac{1}{\sqrt{2}}\).
\[
A^{-1}=\frac{1}{1-2r^{2}}
\begin{bmatrix}1&-r\\-2r&1\end{bmatrix}.
\]
\end{tcolorbox}

\begin{tcolorbox}[title=Key Formulas and Facts, colframe=black!60!white, colback=gray!2]
\begin{itemize}
\item \textbf{Matrix Invertibility Tests:}
  \begin{itemize}
  \item $\det(A)\neq0 \iff A$ invertible
  \item $A$ invertible $\iff$ $\mathrm{rank}(A)=n \iff$ rref$(A)=I$
  \end{itemize}
\item \textbf{Inverse of $2\times2$:}\quad
$A^{-1}=\dfrac{1}{ad-bc}\begin{bmatrix}d&-b\\-c&a\end{bmatrix}$
\item \textbf{Linear Transformation:}\;
$T(\mathbf{x})=A\mathbf{x}$, \;
$T$ linear $\iff$ $T(c\mathbf{x}+\mathbf{y})=cT(\mathbf{x})+T(\mathbf{y})$
\item \textbf{Composition:}\; $T_B\!\circ T_A = T_{BA}$
\item \textbf{Projection onto a line:}\;
$\mathrm{proj}_{\mathbf{u}}(\mathbf{x})=
\dfrac{\mathbf{x}\cdot\mathbf{u}}{\mathbf{u}\cdot\mathbf{u}}\mathbf{u}$
\item \textbf{Reflection across a line:}\;
$\mathrm{ref}_{L}(\mathbf{x})=2\,\mathrm{proj}_{L}(\mathbf{x})-\mathbf{x}$
\item \textbf{Rotation matrix:}\;
$\begin{bmatrix}\cos\theta&-\sin\theta\\\sin\theta&\cos\theta\end{bmatrix}$ preserves length.
\item \textbf{General solution of $A\mathbf{x}=\mathbf{b}$:}\;
$\mathbf{x}=\mathbf{x}_p+\mathbf{x}_h$
\end{itemize}
\end{tcolorbox}
% ---------------------- Quick Reference -- Matrix Operations, Rank, Determinant & Solving A x = b (compact) ----------------------
\begin{tcolorbox}[title={Quick Reference -- Matrix Operations, Rank, Determinant \& Solving $A\vec{x}=\vec{b}$}, colframe=black!60!white, colback=gray!2, fontupper=\tiny]
\begingroup\setlength{\arraycolsep}{1.2pt}\renewcommand{\arraystretch}{0.9}
\textbf{Matrix operations.}
\[
A=\begin{bmatrix}1&2\\3&4\end{bmatrix},\;
B=\begin{bmatrix}5&6\\7&8\end{bmatrix},\quad
A+B=\begin{bmatrix}6&8\\10&12\end{bmatrix},\;
A-B=\begin{bmatrix}-4&-4\\-4&-4\end{bmatrix},\;
2A=\begin{bmatrix}2&4\\6&8\end{bmatrix}.
\]
\textbf{Multiplication (row$\times$col):}\;
\(
AB=\begin{bmatrix}
1\cdot5+2\cdot7 & 1\cdot6+2\cdot8\\
3\cdot5+4\cdot7 & 3\cdot6+4\cdot8
\end{bmatrix}
=\begin{bmatrix}19&22\\43&50\end{bmatrix}
\).
\;Size rule: \((m\times n)(n\times p)\to(m\times p)\).

\textbf{REF vs.\ RREF examples.}
\[
\text{REF }
\left[\!\begin{array}{ccc|c}
1&2&-1&3\\
0&1&4&5\\
0&0&0&1
\end{array}\!\right]
\quad
\text{RREF }
\left[\!\begin{array}{ccc|c}
1&0&-1&1\\
0&1&2&2\\
0&0&1&-1
\end{array}\!\right]
\]
REF: zeros below pivots. RREF: pivots are 1 and the only nonzero in their columns.

\textbf{Rank and nullity.}\;
\(\mathrm{rank}(A)=\#\) pivot cols,\;
\(\text{nullity}(A)=n-\mathrm{rank}(A)\).
If \(\mathrm{rank}([A\mid \vec b])>\mathrm{rank}(A)\Rightarrow\) no solution.
If consistent: free vars \(\Rightarrow\) infinitely many; no free vars \(\Rightarrow\) unique.

\textbf{Solving \(A\vec{x}=\vec{b}\).}\;
1) Form \([A\mid \vec b]\).\;
2) Row-reduce to RREF.\;
3) Pivots \(\Rightarrow\) solved vars; nonpivots \(\Rightarrow\) parameters.\;
4) \(\vec x=\vec x_p+\vec x_h\) (particular + homogeneous).

\textbf{Determinant essentials.}\;
Row swap \(\Rightarrow\) sign flip;\;
\(R_i\gets kR_i\Rightarrow\) det\(\times k\);\;
\(R_i\gets R_i+kR_j\Rightarrow\) det unchanged;\;
triangular \(U\): \(\det(U)=\prod u_{ii}\);\;
square \(A\): \(\det(A)\neq0\Leftrightarrow \mathrm{rank}(A)=n\Leftrightarrow \text{rref}(A)=I_n\) (invertible).

\textbf{Examples.}\;
\(\det\!\begin{bmatrix}a&b\\ c&d\end{bmatrix}=ad-bc\),\quad
\(\det\!\begin{bmatrix}1&2&3\\ 4&5&6\\ 7&8&9\end{bmatrix}=0\Rightarrow \mathrm{rank}=2<3\) (not invertible).
\endgroup
\end{tcolorbox}
% ---------------------- How to Solve A x = b (compact guide) ----------------------
\begin{tcolorbox}[title={How to Solve $A\vec{x}=\vec{b}$ (Step-by-step)}, colframe=black!60!white, colback=gray!2]
\scriptsize
\textbf{Goal.} Find all $\vec{x}$ with $A\vec{x}=\vec{b}$.

\textbf{Recipe.}
\begin{enumerate}[leftmargin=1.1em,itemsep=1pt]
\item Form the augmented matrix $[A\mid \vec{b}]$.
\item Row-reduce $[A\mid \vec{b}]$ to \textbf{RREF} using elementary row ops (swap, scale a row, add multiple of a row).
\item \textbf{Check consistency:} any row $[0\ 0\ \cdots\ 0\mid c]$ with $c\neq0$ means \emph{no solution}.
\item Identify \textbf{pivot columns} (leading 1’s) and \textbf{free} columns (no pivot).
\item Solve pivot variables in terms of the free variables (parameters $r,s,\ldots$).
\item Write the general solution as
\[
\vec{x}=\vec{x}_p+\vec{x}_h,\qquad
\vec{x}_h=\text{all solutions to }A\vec{x}=\vec{0}.
\]
\end{enumerate}

\textbf{Rank test (useful shortcut).}
\[
\mathrm{rank}([A\mid\vec{b}])>\mathrm{rank}(A)\ \Rightarrow\ \text{no solution}.
\]
If consistent:
\[
\text{free vars present }\Rightarrow\ \text{infinitely many solutions},\quad
\text{no free vars }\Rightarrow\ \text{unique solution}.
\]
For square $n\times n$: unique solution $\Leftrightarrow \mathrm{rank}(A)=n$.

\textbf{Mini example (free variable).}
\[
\left[
\begin{array}{ccc|c}
1&2&1&4\\
0&1&3&2
\end{array}
\right]
\ \longrightarrow\ 
\left[
\begin{array}{ccc|c}
1&0&-5&0\\
0&1&3&2
\end{array}
\right].
\]
Let $x_3=t$ (free). Then $x_2=2-3t$, $x_1=5t$.
\[
\vec{x}=\begin{bmatrix}0\\2\\0\end{bmatrix}
+t\begin{bmatrix}5\\-3\\1\end{bmatrix}
=\ \underbrace{\begin{bmatrix}0\\2\\0\end{bmatrix}}_{\vec{x}_p}
+\ t\ \underbrace{\begin{bmatrix}5\\-3\\1\end{bmatrix}}_{\text{one basis vector of }\vec{x}_h}.
\]

\textbf{Mini example (no solution).}
\[
\left[
\begin{array}{ccc|c}
1&2&-1&3\\
0&1&4&5\\
0&0&0&1
\end{array}
\right]
\ \Rightarrow\ \text{row }[0\ 0\ 0\mid 1]\ \Rightarrow\ \text{inconsistent (no solution)}.
\]
\end{tcolorbox}
% ---------------------- Problem 2 and 3 — Explanations and Solutions (compact) ----------------------
\begin{tcolorbox}[title=Problems 2 and 3 — Explanations and Solutions, fontupper=\scriptsize]
\noindent
\textbf{Problem 2.}\; Given the RREF
\[
\begin{bmatrix}
1&0&0&1&0\\
0&1&0&2&1\\
0&0&1&3&2\\
0&0&0&0&0
\end{bmatrix},
\]
pivots are in \(x_1,x_2,x_3\); \(x_4\) is free. Rows give
\[
x_1+x_4=0,\quad x_2+2x_4=1,\quad x_3+3x_4=2.
\]
Let \(x_4=r\). Then
\[
x_1=-r,\quad x_2=1-2r,\quad x_3=2-3r,
\quad
\vec{x}=
\begin{bmatrix}x_1\\x_2\\x_3\\x_4\end{bmatrix}
=
\begin{bmatrix}0\\1\\2\\0\end{bmatrix}
+r
\begin{bmatrix}-1\\-2\\-3\\1\end{bmatrix}.
\]

\medskip
\textbf{Problem 3.}\; System
\[
\begin{cases}
x_1-x_3=1,\\
2x_1-2x_2=1,\\
3x_1-4x_2+x_3=a
\end{cases}
\quad\Longrightarrow\quad
\left[
\begin{array}{ccc|c}
1&0&-1&1\\
2&-2&0&1\\
3&-4&1&a
\end{array}
\right]
\to
\left[
\begin{array}{ccc|c}
1&0&-1&1\\
0&1&-1&\tfrac12\\
0&0&0&\,a-1
\end{array}
\right].
\]
Last row: \(0= a-1\).
\[
a\neq 1 \Rightarrow \text{no solution};\qquad
a=1 \Rightarrow \text{consistent with one free var} \Rightarrow \text{infinitely many};
\]
never exactly one (missing a pivot). For \(a=1\):
\[
x_1-x_3=1,\quad x_2-x_3=\tfrac12,\quad
x_3=r \Rightarrow x_1=1+r,\ x_2=\tfrac12+r,
\]
\[
\vec{x}=
\begin{bmatrix}x_1\\x_2\\x_3\end{bmatrix}
=
\begin{bmatrix}1\\[2pt]\tfrac12\\[2pt]0\end{bmatrix}
+r
\begin{bmatrix}1\\1\\1\end{bmatrix}.
\]
\textit{Summary:}\; (a) \(a\neq1\): no solution;\;
(b) never exactly one;\;
(c) \(a=1\): infinitely many (free var \(x_3\)).
\end{tcolorbox}


\end{multicols}
\end{document}

